\documentclass[12pt]{article}
\usepackage[utf8]{inputenc}
\usepackage[T1]{fontenc}
\usepackage{hyperref}
\title{Yambeta Language Resources for K-12 Curriculum Development}
\author{Compiled by Claude Sonnet 5.5 (claude-sonnet-5-5) for Andr\'as Kornai}
\date{}
\begin{document}
\maketitle

\section{Introduction}
Yambeta (Yambetta, Njambeta, Nedek, Nigi) is a Southern Bantoid language of the Mbam-et-Inoubou department in central Cameroon, around Kon-Yambetta. Speaker estimates vary: 3,700 (Wikipedia, from 1982), about 6,000 (Omniglot, Webonary) and 13,000 for the Nedek people (Joshua Project). Sources name two or four dialects, with Nedek and Nigii as the main ones. Omniglot calls it endangered because few young people learn it. Joshua Project says it is taught in some primary schools.

\section{Primary Educational Resources}
No textbooks or primers were found. Joshua Project reports that the language is taught in some primary schools: \url{https://joshuaproject.net/people_groups/15973}. No curriculum details were found.

\section{Linguistic and Reference Materials}
\begin{itemize}
\item Wikipedia article, which includes phonology: \url{https://en.wikipedia.org/wiki/Yambeta_language}
\item Omniglot page on the language and its Latin-based writing: \url{https://omniglot.com/writing/yambeta.htm}
\item Webonary site with an introduction (and apparently a dictionary): \url{https://www.webonary.work/yambeta/overview/introduction/}
\item SIL Cameroon archive search results for yat: \url{https://www.silcam.org/fr/resources/search/contributors/27378/language/yat/subject/339411}
\end{itemize}

\section{Religious and Community Resources}
Scripture Earth lists scripture materials, including the Gospel of Luke: \url{https://scriptureearth.org/00eng.php?iso=yat}. Joshua Project language page: \url{https://joshuaproject.net/languages/yat}.

\section{Digital and Audio Resources}
No audio or video resources were found beyond the scripture platforms mentioned above.

\section{Curriculum Development Context}
Yambeta is used in school as a subject in some schools, according to the reference description and Joshua Project. An orthography exists, which supports literacy work. The Webonary site and the SIL Cameroon materials are the likeliest sources of vocabulary and reading texts. Details on the schools involved were not found and would have to be established locally.

\end{document}